\section*{Заключение}
\addcontentsline{toc}{section}{Заключение}

В настоящей работе разработан и реализован модуль автоматической генерации
персонализированных учебных конспектов по математике ЕГЭ, центральным
компонентом которого является пайплайн построения корректных
математических иллюстраций по тексту. В данном заключении полученные
результаты соотнесены с целью и задачами, сформулированными во введении,
обозначены теоретическая и практическая значимость работы, а также
зафиксированы её ограничения и направления дальнейшего развития.

\paragraph{Соответствие достигнутых результатов поставленным задачам.}
Цель работы --- разработать и реализовать модуль генерации
персонализированных учебных конспектов, включающих корректные визуальные
материалы, --- достигнута. Каждая из пяти задач, поставленных во введении,
завершена в объёме, достаточном для прототипной демонстрации работоспособности
контура и для перехода к стадии практического внедрения.

Первая задача --- анализ существующих подходов к персонализации учебного
контента, генерации учебных материалов и автоматическому построению
математических визуализаций --- решена в рамках обзора литературы и раздела
об архитектуре системы. Показано, что среди существующих образовательных
платформ ни одна не закрывает полный цикл от аналитики компетенций ученика
к персонализированному конспекту с математическими иллюстрациями: одни
системы реализуют адаптивный подбор задач без генерации объяснительного
контента, другие работают в диалоговом формате без структурированных
визуальных материалов. Показано также, что прямая генерация изображений
моделями текста к изображению (text\nobreakdash-to\nobreakdash-image)
систематически нарушает точные геометрические отношения и не обеспечивает
корректного отображения математических подписей и формул. Это определило
выбор двухшагового процесса с промежуточным структурированным
представлением сцены.

Вторая задача --- проектирование архитектуры процесса генерации конспектов
от вектора компетенций ученика к структурированному тексту с визуальными
элементами --- завершена и закреплена в виде программного контура,
описанного в подразделе~\ref{sec:ch3:architecture}. Контур разбит на
четыре ступени с фиксированными контрактами по входу и выходу: разбиение
конспекта на смысловые секции, двухпроходный планировщик визуалов,
конвертер педагогической таксономии в технические сцены и движок
векторного рендеринга. Такая декомпозиция обеспечивает локализуемость
ошибок и независимое развитие ступеней, что подтверждено практикой
итеративного улучшения прототипа (подраздел~\ref{sec:ch3:improvements}).

Третья задача --- разработка схемы JSON\nobreakdash-представления для
математических визуализаций и реализация движка SVG\nobreakdash-рендеринга
на её основе --- решена в составе модулей \texttt{pipeline/}. Введён
универсальный формат сцены, поддерживающий три семейства иллюстраций
(\texttt{function\_plot}, \texttt{geometry}, \texttt{diagram}); реализован
геометрический рендерер с солвером ограничений, обрабатывающим типовые
построения школьной геометрии (медианы, высоты, биссектрисы, описанные и
вписанные окружности, перпендикуляры, пересечения прямых, средние линии,
касательные, сечения). Стереометрические сцены поддержаны через косую
кабинетную проекцию с автоматической определёнкой видимых и скрытых рёбер
по списку граней. Корректность рендеринга на типовых сценариях
продемонстрирована примерами в подразделах~\ref{sec:ch3:integration}
и~\ref{sec:ch3:geometry}, а также набором офлайн\nobreakdash-тестов,
не требующих обращения к внешнему API.

Четвёртая задача --- интеграция модуля визуализаций в общий процесс
генерации конспектов с учётом профиля ученика --- решена в виде сквозного
конвейера от исходного Markdown\nobreakdash-конспекта до набора SVG
с сохранением промежуточных артефактов на каждом этапе
(подраздел~\ref{sec:ch3:architecture}). Дополнительно реализована единая
точка входа \codett{pipeline/run\_unified.py} с подкомандами для конспектной
и геометрической ветвей, что устраняет дублирование флагов и упрощает
автоматизацию смешанных пакетных прогонов
(см. итерацию~13 в табл.~\ref{tab:ch3:iterations}).

Пятая задача --- оценка качества генерируемых конспектов --- решена
в объёме, соответствующем прототипной стадии. Проведены три независимых
эксперимента с фиксированным протоколом и сохранением подробных метрик:
\begin{enumerate}
    \item Серия повторных прогонов на корпусе из четырёх предметных
    конспектов (подраздел~\ref{sec:ch3:integration},
    табл.~\ref{tab:ch3:run-metrics}). Полная серия выявила содержательное
    наблюдение о повышенной дисперсии планировщика на конспекте по
    вероятности и подтвердила устойчивость показателей по токенам и
    времени в диапазоне 11--13~\% относительно среднего.
    \item Серия из трёх повторных прогонов на стабильном корпусе из
    восьми задач планиметрии и стереометрии
    (подраздел~\ref{sec:ch3:limitations},
    табл.~\ref{tab:ch3:geometry-variance}). Серия зафиксировала среднюю
    задержку полного прогона около 17~с и нулевой fix\nobreakdash-rate
    по двадцати четырём ячейкам; на объединённой выборке из тридцати
    шести ячеек двух серий доля прогонов с инициированным повторным
    запросом составила приблизительно 0{,}028.
    \item Сравнительный эксперимент двух языковых моделей разной
    размерности (\codett{Llama-3.3-70B-Instruct-Turbo} и
    \codett{Qwen2.5-7B-Instruct-Turbo}) на одном корпусе при одинаковом
    протоколе (табл.~\ref{tab:ch3:model-comparison}). Эксперимент
    количественно подтвердил, что бо\'{}льшая модель обеспечивает нулевой
    fix\nobreakdash-rate, тогда как меньшая модель показала
    fix\nobreakdash-rate около 17~\% и одну ячейку отказа из двадцати
    четырёх.
\end{enumerate}
Дополнительно эмпирически подтверждён эффект промоутирования
пост\nobreakdash-валидатора визуализации ответа из режима диагностического
предупреждения в режим замкнутого цикла исправления
(табл.~\ref{tab:ch3:answer-vis-before-after}): доля ячеек с финальным
замечанием снизилась с 5/9 до 0/9 при стоимости пяти дополнительных
вызовов модели на девять прогонов и без структурных регрессов на принятых
кандидатах.

\paragraph{Теоретическая и методологическая значимость.}
В работе систематизированы и обобщены аргументы в пользу двухшаговой
схемы генерации математических иллюстраций с промежуточным
структурированным представлением. Принятая декомпозиция, в которой
языковая модель отвечает за семантику задачи, а специализированный движок
--- за геометрические и арифметические аспекты построения, согласуется с
общей закономерностью, сформулированной по итогам тринадцати
итеративных улучшений прототипа: устойчивость контура повышается тем
сильнее, чем больше шагов рассуждения переносится с языковой модели на
детерминированные алгоритмы. Эта закономерность имеет общий характер и
применима к задачам структурированной генерации за пределами
рассмотренной предметной области (например, к генерации диаграмм в
химии, биологии, экономике), где требуется одновременно интерпретация
естественноязыковой постановки и точное построение по фиксированным
правилам.

Дополнительный методологический вклад связан с введением замкнутого
цикла программной верификации педагогического качества плана. Предложенная
схема, в которой пост\nobreakdash-валидатор по типу ожидаемого ответа
(длина, высота, расстояние, угол, площадь, радиус, периметр, объём,
доказательство, построение) проверяет наличие соответствующего
геометрического объекта и инициирует адресный fix\nobreakdash-запрос
с защитой от тихих регрессов, представляет собой переносимый шаблон
проектирования: предметная сторона валидации меняется, общая структура
(детектор класса ответа, набор требований, контролируемый цикл с условием
строгого улучшения) сохраняется. Соответствующий код реализован в
функциях \codett{\_validate\_answer\_visualization} и
\codett{\_build\_answer\_vis\_fix\_request} модуля \codett{geometry\_auto.py}
и покрыт офлайновыми тестами.

\paragraph{Практическая значимость.}
Реализованный модуль встраивается в действующую платформу подготовки
к ЕГЭ <<Прорыв>>~\cite{Proriv} в качестве компонента, отвечающего за
автоматическое формирование персонализированных учебных материалов.
Автоматизация полного цикла «текст \(\to\) план \(\to\) сцена \(\to\) SVG»
устраняет необходимость ручного создания иллюстраций для каждого
индивидуального профиля компетенций и делает практически достижимой
генерацию учебных материалов под конкретного ученика в близком к реальному
времени. Сохранение промежуточных артефактов в виде самостоятельных
JSON\nobreakdash-файлов на каждом этапе обеспечивает воспроизводимость
прогонов и снижает стоимость диагностики типовых сбоев на порядок:
характерные расследования сводятся к проверке одного промежуточного файла
без повторных вызовов модели. Для геометрической ветви эта же практика
позволяет отделить процесс отладки рендеринга от стоимости обращения к
языковой модели и проводить инженерные эксперименты на сохранённых планах.

Архитектурно предложенный подход применим и за пределами математики ЕГЭ.
Промежуточный JSON\nobreakdash-формат и принцип разделения ответственности
между планировщиком, конвертером и движком переносимы на любую предметную
область, где требуется структурированная генерация графических иллюстраций
под индивидуальный учебный контент.

\paragraph{Ограничения работы.}
Полученные количественные оценки следует читать как точечные выборочные
оценки на прототипном корпусе, а не как окончательные показатели качества.
Все интеграционные эксперименты проведены на корпусе ограниченного размера
(четыре предметных конспекта и восемь геометрических задач), а
формальная сравнительная оценка с альтернативной базовой линией ---
прямой генерацией SVG моделью текста к изображению --- в настоящей
редакции не проведена. Используемые в прототипе показатели имеют
преимущественно технический характер (доля прогонов без ошибок, средняя
задержка, число токенов, fix\nobreakdash-rate), тогда как педагогическое
качество иллюстраций требует независимой экспертной разметки силами
учителей математики. Дополнительные ограничения связаны с двойственностью
текущей реализации рендеринга: сосуществование полноценного сервиса
\texttt{svg-generator} и встроенного запасного рендерера в
\texttt{run\_pipeline.py} обусловлено инфраструктурными причинами на
конкретной рабочей машине и в продуктовой версии должно быть устранено.
Перечисленные ограничения подробно зафиксированы в
подразделе~\ref{sec:ch3:roadmap} и отделены от установленных результатов
работы.

\paragraph{Направления развития.}
По итогам работы намечены следующие задачи ближайшего этапа,
систематизированные в подразделе~\ref{sec:ch3:roadmap}. Первая ---
расширение корпуса до десятков конспектов и сотен геометрических задач
с проведением серии повторных прогонов и формированием статистически
обеспеченных оценок. Вторая --- проведение сравнительного эксперимента
с базовой линией text\nobreakdash-to\nobreakdash-image по согласованным
группам показателей (технические свойства прогона, геометрическая
корректность объектов сцены, педагогическая информативность по протоколу
независимой экспертной разметки). Третья --- введение формального набора
педагогических метрик и согласование протокола экспертной разметки с
руководителем работы. Четвёртая --- унификация рендеринга в единое
окружение с устранением необходимости в запасном рендерере. Пятая ---
расширение типологии визуалов в направлении векторных операций,
анимированных стереометрических построений и сцен с подвижными точками.
Шестая --- объединение конспектной и геометрической ветвей конвейера через
делегирование пошаговой геометрии из планировщика конспекта при
обнаружении соответствующих задач. Седьмая --- систематический эксперимент
с моделями разной размерности и архитектуры (включая рассуждающие модели)
для количественного разнесения вклада размера модели, способа обращения
к ней и структуры самого контура.

\paragraph{Итог.}
Поставленная во введении цель достигнута: разработан и реализован
прототип модуля автоматической генерации персонализированных учебных
конспектов с математическими иллюстрациями, обеспечивающий геометрическую
корректность визуальных материалов и встраиваемый в действующую
образовательную платформу. Результаты прототипа подтверждают
работоспособность принятого архитектурного решения с двухшаговой схемой
и промежуточным структурированным представлением сцены.
Зафиксированные ограничения переведены в перечень задач ближайшего этапа,
выполнение которых превратит полученные качественные наблюдения в
количественные оценки и обеспечит более строгую эмпирическую базу для
заявленных архитектурных выводов.
